\chapter{Results} \label{Results}

The study evaluates three VLMs, Qwen2.5-VL, Gemini 3.1 Pro Preview, and GPT-4o, as judges of text-in-image editing quality, using 623 stimuli that span 13 degradation dimensions and 8 edit types. Each stimulus was evaluated under two experimental conditions: Experiment 1 (perceptual sensitivity, comparing ground truth to degraded image) and Experiment 2 (edit judgment, comparing source to degraded image given an instruction). Together, these yield 3,729 model responses.

\textbf{Experiment 2 is the benchmark-facing task and is the primary analysis.} It replicates the conditions under which VLM judges are deployed in text-editing benchmarks. The research questions are answered from Experiment 2 evidence. Experiment 1 is reported as a supplementary sidecar (Section~\ref{sec:exp1-sidecar}) because, while it documents raw perceptual discrimination capacity, it cannot answer questions about benchmark judgment reliability.

The results are organised around the three research questions, all answered from Experiment 2:

\textbf{RQ1} (\textit{Overall reliability}) is addressed in Section~\ref{sec:score-descriptives}, which characterises each model's aggregate score sensitivity across the full degraded stimulus set, and in Section~\ref{sec:control-conditions}, which establishes the control baselines against which those scores are interpreted.

\textbf{RQ1a} (\textit{Error-type variation}) is addressed in Section~\ref{sec:per-dim-sensitivity}, which breaks down magnitude sensitivity by degradation dimension and model.

\textbf{RQ1b} (\textit{Magnitude transitions}) is addressed in Section~\ref{sec:thresholds}, which estimates the degradation magnitudes at which each model's quality scores first fall below the correct-edit baseline.

The central finding is that VLM judge reliability is both dimension-dependent and model-dependent to a degree that calls the general use of these models as benchmark judges into question. Models show meaningful score sensitivity on image-level degradations (blur, Gaussian noise) but fail to penalise colour, font, and spatial errors until magnitudes far above what constitutes a perceptible flaw. Section~\ref{sec:exp1-sidecar} shows that this pattern is not merely an artefact of the benchmark-task framing: models that cannot see an error in the perceptual task (Experiment 1) do not catch it in the judgment task either. These findings are interpreted further in Chapter~5.


\section{Data Quality and Corpus Overview} \label{sec:data-quality}

\subsection{Corpus Composition}

% Note (from source): not because they are perfect and noop?
The evaluation corpus contains 623 stimuli drawn from a generated set of 703. The difference reflects stimuli that were excluded because they could not be matched across both experiments or because at least one model returned no parseable response for every relevant sub-condition. All analyses operate on this shared corpus of 623 stimuli unless otherwise noted.

The 623 stimuli span 13 degradation dimensions and 8 edit types. Coverage is uneven by design: the stimulus matrix prioritises depth within dimensions over complete factorial balance. Table~\ref{tab:corpus-coverage} shows the total number of stimuli per degradation dimension, collapsed across edit types.

\begin{table}[htbp]
\centering
\caption{Stimulus counts by degradation dimension.}
\label{tab:corpus-coverage}
\begin{tabular}{lc}
\toprule
Degradation dimension & Total stimuli \\
\midrule
\texttt{scale\_error}       & 84 \\
\texttt{font\_weight}       & 79 \\
\texttt{font\_family}       & 76 \\
\texttt{position\_offset}   & 54 \\
\texttt{gaussian\_noise}    & 54 \\
\texttt{opacity}            & 54 \\
\texttt{letter\_spacing}    & 48 \\
\texttt{jpeg\_compression}  & 48 \\
\texttt{color\_offset}      & 44 \\
\texttt{blur}               & 42 \\
\texttt{font\_style}        & 38 \\
\texttt{rotation}           & 33 \\
\bottomrule
\end{tabular}
\end{table}

The coverage table reflects the per-edit-type breakdown described in Section~\ref{Sample Size}. \texttt{color\_offset} is applied exclusively to colour-change stimuli; \texttt{scale\_error} spans seven edit types; and several dimensions (\texttt{gaussian\_noise}, \texttt{blur}, \texttt{jpeg\_compression}) function as cross-cutting image-level degradations applicable to a broad range of edit types.

\subsection{Response Counts and Parse Success Rates}

Each of the 623 stimuli was submitted to each of the three models under both Experiment 1 and Experiment 2, yielding a target of $623 \times 3 \times 2 = 3{,}738$ responses. The achieved total is 3,729, reflecting eight missing responses across conditions due to API timeouts or model refusals that produced no parseable output.

\begin{table}[htbp]
\centering
\caption{Response counts and parse success rates by model and experiment.}
\label{tab:parse-rates}
\begin{tabular}{lcccc}
\toprule
Model & Exp 1 records & Exp 1 parse rate & Exp 2 records & Exp 2 parse rate \\
\midrule
Qwen2.5-VL              & 623 & 100.0\% & 623 & 100.0\% \\
Gemini 3.1 Pro Preview  & 615 & 100.0\% & 623 &  92.3\% \\
GPT-4o                  & 623 & 100.0\% & 622 &  99.8\% \\
\bottomrule
\end{tabular}
\end{table}

Parse success in Experiment 1 is 100\% for all three models. In Experiment 2, Gemini 3.1 Pro Preview shows a parse success rate of 92.3\%, meaning approximately 47 responses were excluded due to format failures, either malformed JSON, truncated outputs, or responses that contained prose rather than the requested structured object. Qwen2.5-VL and GPT-4o achieve near-perfect parse rates in Experiment 2. All analyses involving Gemini Experiment 2 scores are therefore based on 575 records rather than 623, distributed unevenly across dimensions.

% Note (from source): Just rerun the missing ones?
Parse failures are not randomly distributed across stimuli: they concentrate in Gemini's responses to stimuli with complex instructions. This introduces a small selection bias, Gemini's Experiment 2 results are drawn from the subset of stimuli for which the model successfully completed the response format, which may be easier or more unambiguous stimuli. This is noted as a caveat wherever Gemini Experiment 2 data is reported.

\subsection{Stimulus Characteristics}

The per-model corpus used in Experiment 1 closely mirrors the full 623-stimulus corpus. Per-cell counts vary by dimension, for instance, \texttt{rotation} has only 87 stimuli compared to 233 for \texttt{scale\_error}, and within each dimension the number of stimuli per magnitude tier is typically 14--22 for dimensions with three tiers and 4--6 for dimensions with ten tiers (e.g., \texttt{color\_offset}). These counts are sufficient to estimate mean detection rates and detect strong monotone trends, but the statistical power to detect subtle slope differences between models within a single tier is limited for dimensions with small tier counts.


\section{Phase 2 Control Condition Results} \label{sec:control-conditions}

The two Phase 2 control conditions establish the baselines against which all Experiment 2 results are interpreted: the instruction-not-followed control tests whether the rubric is measuring instruction adherence, and the correct-edit control establishes the specificity floor below which degraded-stimulus scores must fall to carry quality-discrimination signal.

\subsection{Instruction-Not-Followed Control: Rubric Validity} \label{sec:noop-control}

The instruction-not-followed control presents the model with the original unedited image as the ``edit result'' alongside the full edit instruction, so no edit was applied. Per-dimension expectations are not uniform. \texttt{instruction\_following} should be 1, since the requested operation was not performed. \texttt{text\_accuracy}, \texttt{visual\_consistency}, and \texttt{layout\_preservation} should be high: the unedited image contains correctly rendered text, visually intact rendering, and an unchanged layout. In the noop condition, the image is unmodified, so the properties these dimensions describe, text correctness, visual integrity, and layout, are intact by construction. The behaviour of \texttt{overall\_quality} is the open question, since the prompt does not define what scope the dimension should encompass, and the label admits multiple plausible readings.

\begin{table}[htbp]
\centering
\caption{Instruction-not-followed control: mean Experiment 2 scores when no edit was applied ($n = 33$--34 per model).}
\label{tab:noop-control}
\begin{tabular}{lccc}
\toprule
Dimension & Qwen2.5-VL & Gemini 3.1 Pro & GPT-4o \\
\midrule
\texttt{instruction\_following} & 1.71 & 1.00 & 1.71 \\
\texttt{text\_accuracy}         & 4.65 & 5.00 & 5.00 \\
\texttt{visual\_consistency}    & 3.71 & 5.00 & 5.00 \\
\texttt{layout\_preservation}   & 5.00 & 5.00 & 5.00 \\
\texttt{overall\_quality}       & 3.62 & 1.00 & 2.94 \\
\bottomrule
\end{tabular}
\end{table}

The four narrow dimensions behave largely as expected. All three models score \texttt{layout\_preservation} at the ceiling, and Gemini and GPT-4o score \texttt{text\_accuracy} and \texttt{visual\_consistency} at the ceiling as well. Qwen's slightly lower scores on \texttt{text\_accuracy} (4.65) and \texttt{visual\_consistency} (3.71) suggest some sensitivity to image properties unrelated to the edit, but the values remain high. \texttt{instruction\_following} is correctly identified as failed by all three models, with Gemini at the floor (1.00) and Qwen and GPT-4o close to it (1.71).

The divergence is concentrated in \texttt{overall\_quality}. Gemini assigns a mean of 1.00, with every noop stimulus receiving the minimum score uniformly across edit types. GPT-4o assigns a mean of 2.94, with 67.7\% of noop stimuli scoring 3 or higher. Qwen assigns a mean of 3.62, with 73.5\% of stimuli scoring 3 or higher. The same stimuli, evaluated against the same prompt, produce scores spanning 2.62 points on a five-point scale depending on which model is judging.

The reading is that \texttt{overall\_quality} is under-specified by the prompt, and that each model resolves the scope according to a different prior. The implications for downstream analysis and benchmark design are taken up in Section~\ref{sec:comparability}; but generally, two consequences follow. First, \texttt{overall\_quality} is reported alongside, rather than in place of, the narrow dimensions throughout Sections~\ref{sec:main-results} and~\ref{sec:exp1-sidecar}, so that readers can see whether a finding holds where the rubric is well-defined. Second, the narrow-dimension mean (\texttt{narrow\_mean}) is used as the primary aggregate where a single Experiment 2 score is needed, since its noop spread across models is 0.41 points compared to 2.62 for \texttt{overall\_quality}.

\subsection{Correct-Edit Control: Specificity Floor (Experiment 2)} \label{sec:correct-edit-control}

% Note (from source): should we instead calculate the baseline from just IF dimension? Is comparing to the ceiling even appropriate, really?
The correct-edit control presents the model with a correctly executed edit as the edit result alongside the instruction that produced it. Because the edit was applied without any degradation (numeric\_magnitude $= 0$), all five rubric dimensions should score at ceiling. A model that fails to reach ceiling on these stimuli is miscalibrated in the direction of over-penalisation: it is penalising edits that contain no error.

\begin{table}[htbp]
\centering
\caption{Correct-edit control: mean Experiment 2 scores and below-ceiling rates ($n = 46$ per model). ``Below-ceiling rate'' is the proportion of stimuli that received a score strictly below 5.}
\label{tab:correct-edit-control}
\resizebox{\textwidth}{!}{%
\begin{tabular}{lcccccc}
\toprule
Dimension & Qwen mean & Qwen $<$5 & Gemini mean & Gemini $<$5 & GPT-4o mean & GPT-4o $<$5 \\
\midrule
\texttt{instruction\_following} & 2.46 & 80.4\% & 3.89 & 29.5\% & 3.72 & 41.3\% \\
\texttt{text\_accuracy}         & 4.91 &  2.2\% & 5.00 &  0.0\% & 5.00 &  0.0\% \\
\texttt{visual\_consistency}    & 4.11 & 30.4\% & 5.00 &  0.0\% & 4.98 &  2.2\% \\
\texttt{layout\_preservation}   & 4.91 &  2.2\% & 5.00 &  0.0\% & 5.00 &  0.0\% \\
\texttt{overall\_quality}       & 4.17 & 52.2\% & 3.98 & 29.5\% & 4.20 & 43.5\% \\
\texttt{narrow\_mean}           & 4.10 & 80.4\% & 4.72 & 29.5\% & 4.67 & 43.5\% \\
\bottomrule
\end{tabular}%
}
\end{table}

The clearest finding is the universal failure to reach ceiling on \texttt{instruction\_following}. Even when no degradation error is present, Qwen scores \texttt{instruction\_following} at a mean of 2.46 with 80.4\% of stimuli below ceiling; Gemini at 3.89 with 29.5\% below ceiling; GPT-4o at 3.72 with 41.3\% below ceiling. This pattern is not explained by edit difficulty, since the stimuli span the full range of edit types and all were correctly executed: the below-ceiling rates reflect systematic under-scoring of correct edits, not genuine instruction failures.

The three dimensions that describe non-instructed properties (\texttt{text\_accuracy}, \texttt{visual\_consistency}, \texttt{layout\_preservation}) are scored at or very close to ceiling across all three models, consistent with the expectation that a correctly-rendered image preserves these properties by construction. The exception is Qwen's \texttt{visual\_consistency}, which has a 30.4\% below-ceiling rate despite no edit error being present; this matches Qwen's broader pattern of misfiring spatial and visual anomaly detectors observed across the identity and noop controls.

The below-ceiling rates on \texttt{narrow\_mean} (80.4\%, 29.5\%, 43.5\%) are driven almost entirely by \texttt{instruction\_following}, since \texttt{text\_accuracy} and \texttt{layout\_preservation} are near-perfect. This means the \texttt{narrow\_mean} specificity floor is similarly degraded: all three models assign below-ceiling \texttt{narrow\_mean} scores to a substantial fraction of correctly executed edits.

A degraded stimulus can only produce a meaningful signal if its score falls below the score the model would have assigned to the correct edit in the first place. Because Qwen assigns a mean \texttt{instruction\_following} of only 2.46 even to perfect edits, its \texttt{instruction\_following} scores on degraded stimuli carry limited diagnostic information: the model cannot distinguish ``edit applied correctly'' from ``edit applied with a small error'' in the \texttt{instruction\_following} dimension.


\section{Main Experiment 2 Results} \label{sec:main-results}

Experiment 2 replicates the actual benchmark task: the model is given the source image, the edit instruction, and the degraded result, and asked to score the result on the five-dimension rubric used by TextEditBench. No ground-truth reference is provided. The model must infer what a correct result would look like from the instruction alone and compare it against the degraded output. This section reports rubric score descriptives, per-dimension magnitude sensitivity, and below-baseline quality-discrimination thresholds. Rubric internal consistency (Cronbach's alpha and intercorrelations) is reported in Section~\ref{sec:supporting-analyses}.

\subsection{Cross-Model Comparability and the Choice of Aggregate} \label{sec:comparability}

Cross-model comparisons and model rankings throughout this section use \texttt{narrow\_mean}, the mean of the four narrowly-scoped rubric dimensions: \texttt{instruction\_following}, \texttt{text\_accuracy}, \texttt{visual\_consistency}, and \texttt{layout\_preservation}. Raw \texttt{overall\_quality} is excluded from cross-model level comparisons. This subsection establishes why.

The instruction-not-followed control in Section~\ref{sec:noop-control} places Gemini at OQ $= 1.00$, GPT-4o at OQ $= 2.94$, and Qwen at OQ $= 3.62$ when no edit is applied. The 2.62-point spread on a five-point scale is the direct consequence of an under-specified prompt. The same stimuli, evaluated against the same prompt, produce non-comparable scores depending on which model is judging.

This scope ambiguity is itself a finding. It supplies worked-example evidence for why decomposed rubrics are the right design choice for VLM-based text edit evaluation, and it connects to a broader pattern in the LLM- and VLM-as-judge literature: holistic verdicts diverge across judge models in ways that decomposed verdicts do not. % Note (from source): check citations: Prosa, Sotelo et al. 2025; matched-holistic study, Zhang 2025
The noop control extends this pattern to text editing in images, in a controlled-stimulus setting where the cross-model disagreement cannot be explained by differential perceptual sensitivity, since the stimuli are identical and unedited. TextEditBench and WeEdit both decompose their evaluation rubrics; the noop divergence on \texttt{overall\_quality} provides empirical justification for that design choice.

The narrow dimensions provide the contrast. Each narrow dimension has a defined scope that every model fills similarly, not always at score 1 on noop stimuli, since \texttt{text\_accuracy} and \texttt{layout\_preservation} legitimately remain near ceiling when the text content and layout are intact, but consistently across models. The noop \texttt{narrow\_mean} spread across the three models is 0.41 points (Qwen 3.77, Gemini 4.00, GPT-4o 4.18), compared to 2.62 points for \texttt{overall\_quality}. Defined scopes produce cross-model comparability even when constituent dimensions have legitimate ceiling effects; under-specified scopes do not. The \texttt{narrow\_mean} baseline is high (rather than near 1) because text content is intact and layout is preserved in noop stimuli, but what matters for cross-model comparability is the between-model consistency, not the absolute level.

Two analytical choices follow. First, \texttt{narrow\_mean} is the primary cross-model quality aggregate throughout Sections~\ref{sec:score-descriptives}--\ref{sec:thresholds}. Model rankings, between-model comparisons of absolute levels, and cross-model claims about per-dimension sensitivity are reported on \texttt{narrow\_mean}. Second, \texttt{overall\_quality} is reported alongside \texttt{narrow\_mean} for completeness, its values appear as reference columns in the descriptive and per-dimension tables, but is interpreted only on its own terms, within a single model.

% Note (from source): how much do we even care about overall quality, actually?

The correct-edit control in Section~\ref{sec:correct-edit-control} shows that \texttt{instruction\_following} is systematically under-scored even when no degradation is present: Qwen 2.46, Gemini 3.89, GPT-4o 3.72. This depresses the per-model correct-edit \texttt{narrow\_mean} baselines to 4.10, 4.72, and 4.67 respectively. Cross-model spread on this baseline is 0.62 points, still substantially narrower than the 2.62-point noop spread on \texttt{overall\_quality}, so \texttt{narrow\_mean} remains the appropriate cross-model aggregate. The per-dimension values in Section~\ref{sec:per-dim-sensitivity} are therefore read against each model's correct-edit baseline rather than an absolute ceiling of 5.

\subsection{Score Descriptives and Rubric Behaviour} \label{sec:score-descriptives}

\begin{table}[htbp]
\centering
\caption{Experiment 2 score descriptives per model, across all 623 degraded stimuli. Ceiling rate = proportion scoring 5; floor rate = proportion scoring 1. \texttt{narrow\_mean} is the primary cross-model quality aggregate; \texttt{overall\_quality} is shown for within-model reference only.}
\label{tab:score-descriptives}
\begin{tabular}{llccc}
\toprule
Dimension & Stat & Qwen2.5-VL & Gemini 3.1 Pro & GPT-4o \\
\midrule
\multirow{3}{*}{\texttt{instruction\_following}}
 & Mean         & 2.80  & 3.68  & 4.29  \\
 & Ceiling rate & 23.1\% & 65.6\% & 78.1\% \\
 & Floor rate   & 35.2\% &  9.7\% & 13.9\% \\
\midrule
\multirow{3}{*}{\texttt{text\_accuracy}}
 & Mean         & 4.81  & 4.95  & 4.99  \\
 & Ceiling rate & 95.0\% & 97.6\% & 99.8\% \\
 & Floor rate   &  4.8\% &  0.9\% &  0.2\% \\
\midrule
\multirow{3}{*}{\texttt{visual\_consistency}}
 & Mean         & 4.13  & 3.50  & 4.41  \\
 & Ceiling rate & 64.7\% & 53.4\% & 69.6\% \\
 & Floor rate   & 17.0\% & 19.1\% &  1.3\% \\
\midrule
\multirow{3}{*}{\texttt{layout\_preservation}}
 & Mean         & 4.97  & 4.46  & 4.78  \\
 & Ceiling rate & 99.2\% & 84.7\% & 87.4\% \\
 & Floor rate   &  0.8\% & 11.7\% &  0.6\% \\
\midrule
\textbf{\texttt{narrow\_mean}} & \textbf{Mean} & \textbf{4.18} & \textbf{4.25} & \textbf{4.62} \\
\midrule
\multirow{3}{*}{\texttt{overall\_quality} (ref.)}
 & Mean         & 4.05  & 3.40  & 4.27  \\
 & Ceiling rate & 40.3\% & 45.4\% & 61.0\% \\
 & Floor rate   &  9.8\% & 19.7\% &  1.0\% \\
\bottomrule
\end{tabular}

\smallskip
\footnotesize\textit{\texttt{narrow\_mean} = mean(instruction\_following, text\_accuracy, visual\_consistency, layout\_preservation).}
\end{table}

Several patterns stand out immediately.

\textbf{\texttt{text\_accuracy} is at ceiling for all models.} Mean scores of 4.81--4.99 and ceiling rates of 95--99.8\% indicate that all three models consistently rate text content accuracy at 5 regardless of the degradation applied. This is expected and appropriate: the degradations in this study do not modify text content, so text accuracy should remain 5 throughout. The consistency of this behaviour across models validates that the \texttt{text\_accuracy} dimension is not confounding the other scores.

\textbf{\texttt{layout\_preservation} shows near-ceiling behaviour overall, with partial model-specific differentiation on spatial degradations.} Global ceiling rates are high (Qwen 99.2\%, GPT-4o 87.4\%, Gemini 84.7\%). Four degradation types do directly modify text element position, orientation, or size, \texttt{position\_offset}, \texttt{rotation}, \texttt{scale\_error}, and \texttt{letter\_spacing}, so the ceiling is not structurally guaranteed. Qwen assigns LP $= 5$ to every stimulus in all four spatial categories (100\% ceiling, $\rho$ n/s throughout); \texttt{layout\_preservation} is fully non-discriminatory for Qwen regardless of magnitude. GPT-4o penalises \texttt{rotation} (LP mean 4.621, 79.3\% ceiling, $\rho = -0.458^*$) and shows significant magnitude correlations for \texttt{letter\_spacing} ($\rho = -0.437^{**}$) and \texttt{scale\_error} ($\rho = -0.291^*$), but treats \texttt{position\_offset} as a non-issue (LP $= 5.000$, 100\% ceiling, $\rho$ n/s). Gemini tracks magnitude for \texttt{position\_offset} ($\rho = -0.451^{**}$) and \texttt{scale\_error} ($\rho = -0.424^{***}$).

\textbf{\texttt{instruction\_following} shows the largest between-model variance, and a primary/secondary split reveals whether that variance reflects genuine instruction sensitivity.} Global IF means are Qwen 2.80, Gemini 4.07, GPT-4o 4.29. To assess whether these scores track actual instruction compliance rather than a general leniency or harshness bias, stimuli are classified as \textit{primary} (the degradation targets the exact property the instruction asked to change, e.g., a \texttt{color\_offset} degradation on a colour-change instruction) or \textit{secondary} (the instruction target was met, but an unrelated property degraded). A well-calibrated model should score IF lower on primary than on secondary; the gap (secondary mean minus primary mean) quantifies this.

Gemini and GPT-4o both show positive gaps (Gemini $+0.56$, GPT-4o $+0.49$), indicating that both models do penalise instruction failure more on primary degradations. Gemini shows especially strong separation on font weight: primary IF of 1.67 (heavy) and 1.80 (light) versus secondary means of 4.60 and 4.58, indicating it reliably assigns near-floor scores when a weight instruction was not executed and near-ceiling scores when a weight instruction was met but something collateral degraded. GPT-4o's gap is similar in direction but attenuated, consistent with its broader leniency and 78.1\% ceiling rate on IF.

Qwen's gap is $-0.41$, in the wrong direction: primary degradations receive higher IF scores than secondary. Qwen's low global mean (2.80) does not reflect sensitivity to instruction failure; it reflects undifferentiated low scoring. The clearest evidence is primary degradation of \texttt{color\_offset}: Qwen assigns IF $= 5.00$ to every one, the same ceiling score it gives to structurally intact noop stimuli. The noop control (Section~\ref{sec:noop-control}), where Qwen's IF noop mean was 1.71 rather than Gemini's universal 1.00, is consistent with this picture: Qwen's IF scores are not anchored to instruction compliance in either direction.

\textbf{Model ranking on \texttt{narrow\_mean}.} Using the cross-model comparable aggregate, the ranking is GPT-4o (4.62) $>$ Gemini (4.25) $>$ Qwen (4.18). All three means are high because \texttt{text\_accuracy} and \texttt{layout\_preservation} are near ceiling for all models, pulling the aggregate up. The discriminatory signal in \texttt{narrow\_mean} comes primarily from \texttt{instruction\_following} and \texttt{visual\_consistency} as Section~\ref{sec:per-dim-sensitivity} will show.

\subsection{Per-Dimension Magnitude Sensitivity in Experiment 2} \label{sec:per-dim-sensitivity}

Table~\ref{tab:exp2-per-dim} reports per-dimension Experiment 2 results. \texttt{narrow\_mean} (NM) is the mean of the four rubric dimensions excluding \texttt{overall\_quality}; OQ is included as a within-model reference. Spearman correlations are BH-corrected.

\begin{table}[htbp]
\centering
\caption{Experiment 2 \texttt{narrow\_mean} (primary aggregate) and \texttt{overall\_quality} (reference, within-model) means per degradation dimension, and Spearman correlations between numeric magnitude and \texttt{narrow\_mean}. \texttt{narrow\_mean} is the primary cross-model quality metric; \texttt{overall\_quality} is included as a reference column for within-model use. Significance: $^*p < 0.05$, $^{**}p < 0.01$, $^{***}p < 0.001$ (BH-corrected). n/s = not significant.}
\label{tab:exp2-per-dim}
\resizebox{\textwidth}{!}{%
\begin{tabular}{l ccc ccc ccc}
\toprule
 & \multicolumn{3}{c}{Qwen2.5-VL} & \multicolumn{3}{c}{Gemini 3.1 Pro} & \multicolumn{3}{c}{GPT-4o} \\
\cmidrule(lr){2-4} \cmidrule(lr){5-7} \cmidrule(lr){8-10}
Dimension & NM & OQ & $\rho$ (NM) & NM & OQ & $\rho$ (NM) & NM & OQ & $\rho$ (NM) \\
\midrule
\texttt{blur}              & 3.79 & 3.38 & $-0.520^{**}$  & 3.04 & \textbf{1.66} & $-0.618^{***}$ & 3.94 & 3.20 & n/s \\
\texttt{gaussian\_noise}   & 3.63 & 2.83 & $-0.510^{***}$ & 3.29 & \textbf{1.57} & $-0.660^{***}$ & 4.25 & 3.54 & $-0.765^{***}$ \\
\texttt{color\_offset}     & \textbf{5.00} & \textbf{5.00} & n/s & 4.90 & 4.65 & $-0.563^{***}$ & 4.81 & 4.50 & $-0.445^{**}$ \\
\texttt{position\_offset}  & 4.09 & 3.98 & n/s & \textbf{4.68} & 4.32 & $-0.684^{***}$ & 4.88 & 4.69 & n/s \\
\texttt{rotation}          & 4.24 & 4.07 & n/s & 4.41 & 3.72 & $-0.596^{**}$  & 4.55 & 4.31 & $-0.448^{*}$ \\
\texttt{scale\_error}      & 4.18 & 4.19 & n/s & 4.46 & 3.82 & $-0.797^{***}$ & 4.62 & 4.21 & $-0.331^{**}$ \\
\texttt{letter\_spacing}   & 4.36 & 4.25 & n/s & 4.46 & 3.56 & $-0.637^{***}$ & 4.68 & 4.39 & $-0.387^{*}$ \\
\texttt{opacity}           & 4.34 & 4.26 & n/s & 4.49 & 3.88 & $-0.542^{***}$ & 4.65 & 4.43 & n/s \\
\texttt{font\_family}      & 4.17 & 4.28 & n/s & 4.07 & 2.97 & $-0.420^{***}$ & 4.52 & 4.08 & $-0.485^{***}$ \\
\texttt{font\_style}       & 4.22 & 4.41 & n/s & 3.99 & 2.82 & n/s & 4.72 & 4.47 & n/s \\
\texttt{font\_weight\_heavy} & 4.18 & 4.19 & n/s & 4.53 & 3.89 & n/s & 4.74 & 4.48 & n/s \\
\texttt{font\_weight\_light} & 3.94 & 3.96 & n/s & 4.51 & 3.67 & n/s & 4.82 & 4.58 & n/s \\
\texttt{jpeg\_compression} & 4.25 & 3.98 & n/s & 4.64 & 4.16 & $-0.490^{**}$  & \textbf{4.95} & \textbf{4.85} & n/s \\
\midrule
\textit{Correct-edit NM baseline} & \textit{4.10} & --- & --- & \textit{4.72} & --- & --- & \textit{4.67} & --- & --- \\
\bottomrule
\end{tabular}%
}

\smallskip
\footnotesize\textit{NM = \texttt{narrow\_mean}; OQ = \texttt{overall\_quality} (raw, within-model reference). $\rho$ (NM) = Spearman(numeric\_magnitude, \texttt{narrow\_mean}), BH-corrected. Correct-edit NM baseline = per-model \texttt{narrow\_mean} on correctly-executed stimuli (Section~\ref{sec:correct-edit-control}); a degraded-stimulus cell at or above this value provides no quality-discrimination signal relative to the model's own ceiling.}
\end{table}

Table~\ref{tab:exp2-per-dim} cells should be read against each model's correct-edit \texttt{narrow\_mean} baseline (Qwen 4.10, Gemini 4.72, GPT-4o 4.67; Section~\ref{sec:correct-edit-control}). Cells at or above the baseline carry no quality-discrimination signal, and the n/s correlation cells where this holds are not merely magnitude-insensitive but unresponsive to the degradation relative to the model's own ceiling.

Three patterns warrant attention.

\textbf{Qwen is categorically insensitive to edit-layer magnitude.} Only \texttt{blur} and \texttt{gaussian\_noise} produce significant correlations; the other 11 dimensions are n/s. \texttt{color\_offset} is the extreme case: NM $=$ OQ $= 5.00$ across all 40 stimuli and 10 magnitude tiers, with zero variance. Gemini ($\rho = -0.563^{***}$) and GPT-4o ($\rho = -0.445^{**}$) detect the same magnitude axis, so the failure is model-specific rather than a stimulus artifact. Against the 4.10 baseline, only blur (3.79), gaussian\_noise (3.63), and font\_weight\_light (3.94) are meaningfully below; position\_offset (4.09) is marginally below; the remaining nine dimensions sit at or above Qwen's correct-edit ceiling.

\textbf{Image-layer degradations are the only dimensions with consistent cross-model signal}, and also the lowest-scored: Gemini NM 3.04 (blur) / 3.29 (noise), Qwen 3.79 / 3.63, GPT-4o 3.94 / 4.25. NM is markedly higher than OQ here (Gemini OQ 1.66 / 1.57) because \texttt{text\_accuracy} and \texttt{layout\_preservation} are unaffected by blur and noise and pull the average up. NM is therefore a conservative aggregate for image-layer degradations.

\textbf{GPT-4o's JPEG inflation survives aggregation.} NM $= 4.95$ is the highest in its column, exceeds its overall NM of 4.62, and sits 0.28 points above the GPT-4o correct-edit baseline of 4.67, GPT-4o rates JPEG-compressed stimuli higher than it rates correctly-executed edits. Gemini scores the same stimuli at NM $= 4.64$ (below its own 4.72 baseline) with $\rho = -0.490^{**}$, so the inflation is model-specific. Section~\ref{sec:exp1-exp2-correlation} documents the correlation between this pattern and Experiment 1 perceptual scores.

The typography cluster (\texttt{font\_style}, \texttt{font\_weight\_heavy}, \texttt{font\_weight\_light}) is n/s across all three models. This is compatible with either genuine cross-model insensitivity or a non-monotonic magnitude axis for font weight (Section~\ref{Detection and Sensitivity Metrics}); the data do not distinguish them.

\subsection{Below-Baseline Quality-Discrimination Thresholds} \label{sec:thresholds}

The Phase 1 perceptual detection thresholds (Section~\ref{sec:phase1-thresholds}) document where each model transitions from unreliable to reliable perceptual discrimination. The correct-edit baselines from Section~\ref{sec:correct-edit-control} raise a distinct question for the judgment task: at what magnitude does the model penalise the error enough to score it below the level it would assign to a correctly-executed edit? These are different constructs. The analysis below estimates the \textbf{quality-discrimination threshold}: the degradation magnitude at which the fitted Experiment 2 \texttt{narrow\_mean} trend crosses each model's correct-edit NM baseline.

\textbf{Method.} For each (model, dimension) cell where the Spearman correlation between \texttt{numeric\_magnitude} and \texttt{narrow\_mean} is BH-significant in Table~\ref{tab:exp2-per-dim}, a linear regression of $\texttt{narrow\_mean} \sim \texttt{numeric\_magnitude}$ is fitted. The crossing point $x^*$ where the fitted line reaches the per-model correct-edit baseline (Qwen 4.10, Gemini 4.72, GPT-4o 4.67) is the quality-discrimination threshold. Bootstrap 95\% CIs on $x^*$ use 1000 resamples of the per-dimension data with replacement. $x^*$ is mapped to the degradation tier it falls into. For n/s cells, no regression is fitted; the cell is classified as ``never below baseline'' (mean NM at or above the baseline despite no magnitude correlation) or ``always below baseline'' (mean NM below baseline despite non-significant correlation, indicating a flat penalty that does not scale with magnitude). % Note (from source): this classification is inaccurate, rerun stats and add the indistinguishable classification.
The $n = 46$ correct-edit sample is adequate for stable per-model baseline estimates (SE $\approx 0.08$) but does not support per-dimension baseline breakdown, so a single per-model baseline is used throughout.

\begin{table}[htbp]
\centering
\caption{Below-baseline quality-discrimination threshold tier per (model, dimension). Cells show the magnitude tier where the fitted \texttt{narrow\_mean} trend crosses the per-model correct-edit baseline, or the categorical outcome for n/s cells. $x^*$ is the estimated crossing magnitude in dimension-native units; 95\% CI from 1000 bootstrap resamples. Cells marked $\dagger$ have CIs that span zero or very large ranges and should be interpreted with caution.}
\label{tab:thresholds}
\begin{tabular}{p{3cm} p{3.8cm} p{3.8cm} p{3.8cm}}
\toprule
Dimension & Qwen (baseline 4.10) & Gemini (baseline 4.72) & GPT-4o (baseline 4.67) \\
\midrule
\texttt{blur}               & medium ($x^*$=2.55, [1.77, 3.42] rad) & always below baseline & always below baseline \\
\texttt{color\_offset}      & never below baseline & very\_large ($x^*$=27.67, [21.12, 51.40] $\Delta E$) & large ($x^*$=21.11, [16.42, 46.22] $\Delta E$) \\
\texttt{font\_family}       & never below baseline & medium ($x^*$=0.18, [$-$0.39, 0.41] class-dist) $\dagger$ & large ($x^*$=0.58, [0.43, 0.68] class-dist) \\
\texttt{font\_style}        & never below baseline & always below baseline & never below baseline \\
\texttt{font\_weight\_heavy} & never below baseline & always below baseline & never below baseline \\
\texttt{font\_weight\_light} & always below baseline & always below baseline & never below baseline \\
\texttt{gaussian\_noise}    & small ($x^*$=5.49, [$-$10.37, 11.13] $\sigma$) $\dagger$ & always below baseline & small ($x^*$=7.78, [4.55, 10.45] $\sigma$) \\
\texttt{jpeg\_compression}  & never below baseline & medium ($x^*$=37.03, [$-$7.38, 74.04] qual.) $\dagger$ & never below baseline \\
\texttt{letter\_spacing}    & never below baseline & small ($x^*$=2.27, [0.66, 3.46] px) & medium ($x^*$=4.54, [2.34, 8.05] px) \\
\texttt{opacity}            & never below baseline & medium ($x^*$=0.19, [0.07, 0.27] reduction) & always below baseline \\
\texttt{position\_offset}   & always below baseline & medium ($x^*$=34.24, [25.45, 47.63] px) & never below baseline \\
\texttt{rotation}           & never below baseline & tiny ($x^*$=1.43, [$-$25.12, 9.72]$^\circ$) $\dagger$ & medium ($x^*$=6.87, [$-$8.72, 26.10]$^\circ$) $\dagger$ \\
\texttt{scale\_error}       & never below baseline & large ($x^*$=22.73, [19.42, 26.35]\%) & large ($x^*$=25.58, [0.37, 51.83]\%) $\dagger$ \\
\bottomrule
\end{tabular}

\smallskip
\footnotesize\textit{$\dagger$ CI spans zero or covers a very wide range; the crossing-point estimate is not reliable for these cells.}
\end{table}

\textbf{Qwen: almost no discrimination.} The most striking finding in Table~\ref{tab:thresholds} is that Qwen produces ``never below baseline'' for 9 of 13 dimensions. Qwen assigns NM $= 5.00$ to all colour-offset stimuli in Experiment 2 regardless of magnitude: the mean NM sits at the full ceiling, 0.90 points above the Qwen baseline, with zero variance.

Only \texttt{blur} and \texttt{gaussian\_noise} cross the baseline in Experiment 2 for Qwen, consistent with the only two dimensions showing significant NM correlations in Table~\ref{tab:exp2-per-dim}. The \texttt{blur} threshold estimate is reliable ($x^* = 2.55$ rad, medium tier, 95\% CI [1.77, 3.42]).

\textbf{Gemini: early penalisation on image degradations, late on colour.} Gemini's image-layer degradations (blur, Gaussian noise) are classified as ``always below baseline'', the fitted trend is already below the 4.72 baseline at the smallest tested magnitude, indicating that even mild blur and noise depress Gemini's NM below its ceiling level on correct edits. This is consistent with the steep OQ response Gemini shows for these dimensions (Table~\ref{tab:exp2-per-dim}). For spatial and typographic degradations, Gemini has defined thresholds on most dimensions with tight CIs: letter\_spacing (small tier, $x^* \approx 2.3$ px), opacity (medium tier, $x^* \approx 0.19$ reduction), position\_offset (medium tier, $x^* \approx 34$ px), and scale\_error (large tier, $x^* \approx 23\%$). The rotation and jpeg\_compression threshold estimates have very wide CIs and are not interpretable.

\textbf{GPT-4o: variable alignment, with notable flat penalties.} GPT-4o's quality-discrimination threshold for \texttt{gaussian\_noise} is at the small tier ($x^* \approx \sigma$ 7.8, CI [4.6, 10.5]). GPT-4o applies a flat quality penalty for blur regardless of magnitude, which is arguably appropriate, any visible blur should penalise quality, but the lack of gradient means the score does not track severity. \texttt{color\_offset} and \texttt{scale\_error} have wide CIs for GPT-4o and should be interpreted with caution.


\section{Research Question Answers} \label{sec:rq-answers}

This section states direct answers to the three research questions on the basis of the evidence reported in Section~\ref{sec:main-results}.

\subsection*{RQ1: Overall Reliability}

\textbf{RQ1.} \textit{Are VLM judges reliably sensitive to degradation in text-in-image editing outputs?}

% Note (from source): rewrite this section
\textbf{Answer.} Reliability is partial and highly uneven. All three models produce \texttt{narrow\_mean} scores that are statistically above the correct-edit baseline on at least some dimensions (Section~\ref{sec:score-descriptives}), but the cross-model picture is dominated by near-ceiling behaviour. The \texttt{narrow\_mean} aggregate across all 623 degraded stimuli is GPT-4o 4.62, Gemini 4.25, Qwen 4.18, all substantially closer to the correct-edit baselines (4.67, 4.72, 4.10 respectively) than to the floor. The discriminatory signal comes almost entirely from \texttt{instruction\_following} and \texttt{visual\_consistency}. Qwen2.5-VL shows the most severe limitation: colour deviations receive near-maximum scores regardless of magnitude (NM $=$ OQ $= 5.00$ for \texttt{color\_offset}), layout preservation is always 5, and text accuracy is always 5, leaving only \texttt{visual\_consistency} as a potentially discriminatory dimension. Gemini's judgment scoring is more functional as a quality measure: its scores track magnitude more faithfully and its \texttt{instruction\_following} dimension provides real discriminatory power on most error types. GPT-4o occupies an intermediate position with moderate sensitivity on most detected dimensions but low sensitivity on font and spatial position error types.

In sum, VLM judges show some quality-discrimination capacity above their correct-edit baselines on image-level degradations, but the capacity to discriminate is insufficient across the full range of degradation types to constitute reliable evaluation.

\subsection*{RQ1a: Error-Type Variation}

\textbf{RQ1a.} \textit{Does sensitivity vary across error types, and do different VLM judges have the same blind spots?}

\textbf{Answer.} Variation across error types is large and follows a consistent pattern. From the per-dimension analysis in Section~\ref{sec:per-dim-sensitivity}:

\begin{itemize}
    \item \textbf{Image-layer degradations} (blur, Gaussian noise) produce significant magnitude correlations for all three models and are the only dimensions where scores reliably fall below correct-edit baselines across models.
    \item \textbf{Colour offset} is detected by Gemini and GPT-4o (significant NM correlations) but not by Qwen (NM $= 5.00$, zero variance). Even where detected, quality-discrimination thresholds are late (Section~\ref{sec:thresholds}).
    \item \textbf{Spatial and typographic degradations} (font style, font weight, position offset, rotation, scale error, letter spacing) are n/s for Qwen on almost all cells. Gemini shows significant correlations on most spatial and several typographic dimensions; GPT-4o shows significant correlations on font family, letter spacing, rotation, and scale error.
    \item The typography cluster (\texttt{font\_style}, \texttt{font\_weight\_heavy}, \texttt{font\_weight\_light}) is n/s across all three models, indicating universal insensitivity regardless of magnitude.
\end{itemize}

The models therefore do not share the same blind spots. Qwen is categorically insensitive to edit-layer magnitude except for image-level degradations. Gemini and GPT-4o are more sensitive but diverge on specific dimensions (e.g., GPT-4o is sensitive to JPEG compression where Gemini is not). The choice of judge model determines not only which errors are penalised but how reliably each error type is flagged.

\subsection*{RQ1b: Magnitude Transitions}

\textbf{RQ1b.} \textit{At what degradation magnitudes do VLM judges transition from non-penalising to penalising?}

\textbf{Answer.} The quality-discrimination thresholds reported in Section~\ref{sec:thresholds} show that, where crossings exist, they occur at magnitudes that are generally moderate to large, and for many dimensions no crossing is observed at all:

\begin{itemize}
    \item \textbf{Blur and Gaussian noise:} Gemini is ``always below baseline'' (even mild degradation depresses NM below the correct-edit ceiling); GPT-4o applies a flat blur penalty regardless of magnitude. Qwen's blur threshold is at the medium tier (radius $\approx 2.55$ rad); Gaussian noise threshold estimates are unreliable for Qwen.
    \item \textbf{Colour offset:} Gemini's quality-discrimination threshold is at the very\_large tier ($\Delta E \approx 28$). GPT-4o's threshold is at the large tier. Qwen never crosses the baseline.
    \item \textbf{Spatial dimensions:} Where Gemini has defined thresholds, they are at medium (position\_offset $\approx 34$ px) and large (scale\_error $\approx 23\%$) tiers. GPT-4o's crossings for letter spacing and scale error are at medium and large tiers respectively.
    \item \textbf{Typographic font dimensions:} Font style, font weight, and font family produce ``never below baseline'' or ``always below baseline'' without reliable crossing estimates across most (model, dimension) cells.
\end{itemize}

For the majority of error types, current VLM judges do not penalise scores below their correct-edit baseline until magnitudes that substantially exceed the range of errors likely encountered in real text-editing system outputs. This directly limits the validity of VLM-based evaluation in benchmarks that rely on these models as judges.


\section{Supporting Analyses} \label{sec:supporting-analyses}

This section reports two supporting analyses that complement the main Experiment 2 results in Section~\ref{sec:main-results}: rubric internal consistency (Cronbach's alpha and intercorrelation structure) and inter-model agreement on Experiment 2 classifications. Neither analysis answers the research questions directly; they characterise the measurement properties of the rubric and the degree to which model judgments converge.

\subsection{Rubric Internal Consistency and Intercorrelations}

\begin{table}[htbp]
\centering
\caption{Spearman intercorrelations among Experiment 2 score dimensions and Cronbach's alpha, per model. IF = \texttt{instruction\_following}; OQ = \texttt{overall\_quality} (raw, within-model).}
\label{tab:intercorrelations}
\begin{tabular}{lccc}
\toprule
Dimension pair & Qwen2.5-VL $\rho$ & Gemini 3.1 Pro $\rho$ & GPT-4o $\rho$ \\
\midrule
IF $\leftrightarrow$ \texttt{text\_acc}          & 0.26 & 0.23 & 0.08 \\
IF $\leftrightarrow$ \texttt{visual\_cons}       & 0.47 & 0.35 & 0.39 \\
IF $\leftrightarrow$ \texttt{layout\_pres}       & 0.11 & 0.06 & 0.27 \\
IF $\leftrightarrow$ OQ                          & 0.38 & \textbf{0.66} & \textbf{0.77} \\
\texttt{text\_acc} $\leftrightarrow$ \texttt{visual\_cons}  & 0.38 & 0.23 & 0.08 \\
\texttt{text\_acc} $\leftrightarrow$ OQ          & 0.37 & 0.23 & 0.08 \\
\texttt{visual\_cons} $\leftrightarrow$ OQ       & \textbf{0.72} & \textbf{0.79} & \textbf{0.80} \\
\texttt{layout\_pres} $\leftrightarrow$ OQ       & 0.15 & 0.55 & 0.55 \\
\midrule
\textbf{Cronbach's $\alpha$}                     & \textbf{0.762} & \textbf{0.765} & \textbf{0.782} \\
\bottomrule
\end{tabular}
\end{table}

Cronbach's alpha is approximately 0.76--0.78 for all three models, reflecting moderate but not high internal consistency. The rubric dimensions are neither fully redundant nor fully independent. Two structural patterns emerge from the intercorrelation matrix.

First, \texttt{visual\_consistency} is the strongest single predictor of \texttt{overall\_quality} across all models ($\rho = 0.72$--0.80). This is the dimension most directly measuring whether non-target properties of the image changed, that is, whether the degradation affected properties it should not have. The high correlation indicates that \texttt{overall\_quality} is largely driven by \texttt{visual\_consistency}. Which means \texttt{visual\_consistency} carries most of the variance in OQ and including OQ in a composite with the narrow dimensions would double-count it, hence its exclusion from \texttt{narrow\_mean}.

Second, \texttt{text\_accuracy} and \texttt{layout\_preservation} are weakly correlated with \texttt{overall\_quality} for most models. \texttt{layout\_preservation} is almost always rated 5 (Section~\ref{sec:score-descriptives}), so its correlation with \texttt{overall\_quality} is attenuated. \texttt{text\_accuracy} is at ceiling (Section~\ref{sec:score-descriptives}), adding essentially no variance. For Gemini, \texttt{layout\_preservation} correlates more strongly with OQ ($\rho = 0.55$) because Gemini uses the dimension more dynamically.

Third, \texttt{instruction\_following} correlates with OQ most strongly for Gemini and GPT-4o ($\rho = 0.66$ and 0.77 respectively), but more weakly for Qwen ($\rho = 0.38$). This is consistent with the Instruction-Not-Followed control result (Section~\ref{sec:noop-control}): for Gemini and GPT-4o, \texttt{instruction\_following} meaningfully predicts the edit quality assessment; for Qwen, OQ is partly decoupled from instruction adherence.

\subsection{Inter-Model Agreement on Experiment 2 Classifications}

\textit{[WIP --- this section will report Fleiss' kappa and agreement patterns for Experiment 2 quality-discrimination classifications (above/below correct-edit baseline). Note: the current Section~\ref{sec:exp1-sidecar} reports Fleiss' kappa on Phase 1 blind/sensitive classifications; that analysis is based on Experiment 1 detection rates and does not transfer directly to Experiment 2 quality-discrimination. A parallel kappa analysis on Experiment 2 thresholds and/or per-cell sensitivity is pending.]}


\section{Experiment 1: Perceptual Discrimination Sidecar} \label{sec:exp1-sidecar}

Experiment 1 isolated raw visual discrimination: each model was shown the ground-truth image and the degraded image and asked only whether they differ, with no edit instruction and no source image. This task cannot answer the research questions, which concern benchmark judgment reliability, but it provides supplementary context: it documents what the models can see independently of the instruction-reasoning step that Experiment 2 adds. Comparing Experiment 1 and Experiment 2 patterns on the same stimuli shows where the judgment task adds or suppresses signal relative to raw perceptual capacity.

\subsection{Experiment 1 Controls} \label{sec:exp1-controls}

\subsubsection{Identity Control: False Positive Rates}

In the identity control, the model is shown an identical image pair, the correctly edited image presented as both the reference and the comparison, and should report \texttt{detected\_difference: false} and \texttt{similarity\_score: 5}. Any deviation is a false positive. Each control set contains 46 stimuli spanning the edit-type distribution of the main corpus.

False positive rates are heavily concentrated in a single dimension. Excluding \texttt{alignment\_error} (where models appear to fire spatial-anomaly detectors in response to the relocated text's unusual position relative to their internal layout priors, rather than genuinely detecting a difference between identical images), the adjusted model-level FPR collapses to 0.0\% (Qwen), 3.2\% (Gemini), and 3.2\% (GPT-4o). All three models assign mean similarity scores close to 5 on identical pairs (4.96, 4.52, 4.72), confirming that outside of the alignment-error hallucination, the models are effectively unbiased against a null of identical images.

\begin{table}[htbp]
\centering
\caption{Identity control: false positive rates and mean similarity scores (Experiment 1). ``FPR (adjusted)'' excludes \texttt{alignment\_error}.}
\label{tab:identity-control}
\begin{tabular}{lcccc}
\toprule
Model & $n$ & FPR (adjusted) & FPR (all dims) & Mean similarity score \\
\midrule
Qwen2.5-VL              & 46 & \textbf{0.0\%} & 4.3\%  & 4.96 \\
Gemini 3.1 Pro Preview  & 46 &  3.2\%          & 26.1\% & 4.52 \\
GPT-4o                  & 46 &  3.2\%          & 21.7\% & 4.72 \\
\bottomrule
\end{tabular}
\end{table}

The adjusted FPR values (0.0\%, 3.2\%, 3.2\%) serve as each model's individual null baseline for the blind/sensitive classification in Section~\ref{sec:phase1-detection}. A dimension is considered a genuine blind spot only when the model's detection rate on degraded stimuli does not significantly exceed its adjusted FPR (one-sided binomial test with BH correction). Because all adjusted FPRs are near zero, even modest detection rates pass the test.

\subsubsection{Source-Target Control: Sensitivity Floor}

The source-target control pairs the correctly executed edit (as reference) with the original unedited source image (as comparison). This is the maximum-signal condition in Experiment 1: the two images differ by exactly the intended edit, which is the largest difference any stimulus in the set contains. Detection rate on this control therefore establishes an upper bound on the model's per-instance sensitivity. A model that fails here cannot, in principle, reliably detect smaller degradations, and low detection rates on degraded stimuli become uninterpretable, we cannot distinguish a genuine sub-threshold degradation from a model that under-reports differences regardless of magnitude. High detection here is a prerequisite for treating per-dimension sensitivity curves as meaningful.

\begin{table}[htbp]
\centering
\caption{Source-target control: detection rate of the ground-truth edit vs.\ source and mean similarity scores (Experiment 1).}
\label{tab:source-target-control}
\begin{tabular}{lccc}
\toprule
Model & $n$ & Detection rate & Mean similarity score \\
\midrule
Qwen2.5-VL              & 34 & 14.7\%          & 4.74 \\
Gemini 3.1 Pro Preview  & 32 & \textbf{93.8\%} & 3.69 \\
GPT-4o                  & 34 & 55.9\%          & 4.27 \\
\bottomrule
\end{tabular}
\end{table}

Model behaviour diverges sharply. Gemini detects 93.8\% of ground-truth edits and assigns a mean similarity score of 3.69, reflecting genuine sensitivity to the visual difference between source and correctly-edited image. GPT-4o detects 55.9\% with a mean similarity score of 4.27, showing moderate sensitivity. Qwen detects only 14.7\% and gives a mean similarity score of 4.74: for Qwen, the source and the correctly-executed edit look nearly identical, which is consistent with its broader pattern of very low detection rates across most dimensions in the main experiment.

The source-target detection rates establish that the models are ordered Gemini $>$ GPT-4o $>$ Qwen in their sensitivity to actual visual changes. Qwen's low detection rate here confirms that its near-zero detection rates on several degradation dimensions (Section~\ref{sec:phase1-detection}) reflect genuine visual insensitivity to those edit types, not merely task uncertainty.

\subsection{Phase 1 Detection Rates} \label{sec:phase1-detection}

The detection rate on degraded stimuli, interpreted relative to each model's adjusted false positive rate (Section~\ref{sec:exp1-controls}) via BH-corrected binomial tests, classifies each (dimension, model) cell as either \textit{sensitive} (detection rate significantly above FPR) or \textit{blind} (not distinguishable from FPR). Spearman rank correlations between numeric magnitude and similarity score provide a continuous measure of sensitivity that complements the binary classification.

\begin{table}[htbp]
\centering
\caption{Experiment 1 overall detection rates and mean similarity scores across all 623 degraded stimuli.}
\label{tab:exp1-overall}
\begin{tabular}{lccc}
\toprule
Model & Detection rate & Mean similarity score & Similarity score std \\
\midrule
Qwen2.5-VL              & 19.1\% & 4.65 & 0.78 \\
Gemini 3.1 Pro Preview  & 59.0\% & 4.09 & 0.95 \\
GPT-4o                  & 44.1\% & 4.41 & 0.75 \\
\bottomrule
\end{tabular}
\end{table}

In raw terms, Gemini detects the largest proportion of degradations (59.0\%), followed by GPT-4o (44.1\%) and Qwen (19.1\%). All three models have near-zero adjusted FPR, confirming that these detection rates reflect genuine signal.

\begin{table}[htbp]
\centering
\caption{Experiment 1 per-dimension detection rates and magnitude--score correlations. B = blind (BH-corrected $p \geq 0.05$, binomial test vs.\ adjusted FPR); S = sensitive. $\rho$ = Spearman correlation between numeric magnitude and similarity score. Significance: $^*p < 0.05$; $^{***}p < 0.001$ (BH-corrected).}
\label{tab:exp1-per-dim}
\resizebox{\textwidth}{!}{%
\begin{tabular}{l ccc ccc ccc}
\toprule
 & \multicolumn{3}{c}{Qwen2.5-VL} & \multicolumn{3}{c}{Gemini 3.1 Pro} & \multicolumn{3}{c}{GPT-4o} \\
\cmidrule(lr){2-4} \cmidrule(lr){5-7} \cmidrule(lr){8-10}
Dimension & DR & $\rho$ & Class. & DR & $\rho$ & Class. & DR & $\rho$ & Class. \\
\midrule
\texttt{blur}              & 66.7\% & $-0.869^{***}$ & \textbf{S} & 88.1\% & $-0.872^{***}$ & \textbf{S} & 95.2\% & $-0.631^{***}$ & \textbf{S} \\
\texttt{gaussian\_noise}   & 74.1\% & $-0.722^{***}$ & \textbf{S} & 77.8\% & $-0.907^{***}$ & \textbf{S} & 90.7\% & $-0.708^{***}$ & \textbf{S} \\
\texttt{font\_family}      &  5.6\% & $-0.287^{*}$   & \textbf{S} & 70.0\% & $-0.660^{***}$ & \textbf{S} & 54.2\% & $-0.620^{***}$ & \textbf{S} \\
\texttt{font\_style}       &  0.0\% & n/a            & B          & 85.3\% & $-0.200$       & \textbf{S} & 23.5\% & $-0.139$       & \textbf{S} \\
\texttt{font\_weight\_heavy} & 0.0\% & n/a            & B          & 66.0\% & $-0.278$       & \textbf{S} & 27.1\% & $-0.327^{*}$   & \textbf{S} \\
\texttt{font\_weight\_light} & 0.0\% & n/a            & B          & 53.8\% & $-0.022$       & \textbf{S} & 15.4\% & $-0.123$       & \textbf{S} \\
\texttt{letter\_spacing}   & 11.4\% & $-0.501^{***}$ & \textbf{S} & 65.9\% & $-0.644^{***}$ & \textbf{S} & 38.6\% & $-0.656^{***}$ & \textbf{S} \\
\texttt{opacity}           & 24.1\% & $-0.426^{***}$ & \textbf{S} & 52.8\% & $-0.672^{***}$ & \textbf{S} & 38.9\% & $-0.637^{***}$ & \textbf{S} \\
\texttt{jpeg\_compression} & 12.5\% & $-0.179$       & \textbf{S} &  8.7\% & $-0.148$       & B          & 68.8\% & $-0.469^{***}$ & \textbf{S} \\
\texttt{position\_offset}  &  1.9\% & $-0.225$       & \textbf{S} & 41.5\% & $-0.746^{***}$ & \textbf{S} & 13.0\% & $-0.221$       & \textbf{S} \\
\texttt{rotation}          & 24.1\% & $-0.645^{***}$ & \textbf{S} & 48.3\% & $-0.864^{***}$ & \textbf{S} & 34.5\% & $-0.655^{***}$ & \textbf{S} \\
\texttt{scale\_error}      & 14.1\% & $-0.541^{***}$ & \textbf{S} & 62.3\% & $-0.804^{***}$ & \textbf{S} & 29.5\% & $-0.705^{***}$ & \textbf{S} \\
\texttt{color\_offset}     & 10.0\% & $-0.502^{***}$ & \textbf{S} & 40.0\% & $-0.831^{***}$ & \textbf{S} & 27.5\% & $-0.724^{***}$ & \textbf{S} \\
\bottomrule
\end{tabular}%
}
\end{table}

The sensitivity taxonomy emerging from Table~\ref{tab:exp1-per-dim}:

\textbf{Universally sensitive dimensions} (all three models classified sensitive): \texttt{blur}, \texttt{gaussian\_noise}, \texttt{opacity}, \texttt{color\_offset}, \texttt{font\_family}, \texttt{letter\_spacing}, \texttt{position\_offset}, \texttt{rotation}, and \texttt{scale\_error} (9 of 13 dimensions).

\textbf{Partially sensitive dimensions} (two of three models sensitive): \texttt{font\_style}, \texttt{font\_weight\_heavy}, and \texttt{font\_weight\_light} (Gemini and GPT-4o sensitive; Qwen blind); \texttt{jpeg\_compression} (Qwen and GPT-4o sensitive; Gemini blind).

\textbf{True blind spots} --- Qwen on \texttt{font\_style}, \texttt{font\_weight\_heavy}, and \texttt{font\_weight\_light} (0\% detection). Gemini is additionally blind to \texttt{jpeg\_compression}.

The key dimension-level findings that contextualise Experiment 2 patterns: Qwen shows a categorical absence of sensitivity to typographic style and weight changes (0\% detection), and only moderate sensitivity to image-level degradations. Gemini shows the broadest typographic sensitivity (font family 70.0\%, font style 85.3\%, font weight heavy 66.0\%, font weight light 53.8\%) but is blind to JPEG compression. GPT-4o leads on image-level degradations (blur 95.2\%, Gaussian noise 90.7\%) and shows moderate typographic sensitivity.

\subsection{Phase 1 Detection Thresholds} \label{sec:phase1-thresholds}

The threshold analysis uses the discrete JND estimate: the smallest magnitude tier at which detection rate reaches or exceeds 50\%, expressed both as a tier label and as the mean numeric magnitude in physical units for that tier. For dimensions or models where no tier reaches the threshold, the entry is ``NC'' (not crossed).

\begin{table}[htbp]
\centering
\caption{Detection threshold tiers and corresponding numeric magnitudes for the 50\% detection criterion. NC = not crossed (no tier reaches threshold). Units: $\Delta E$ for colour; px for position; \% for scale; degrees for rotation; $\sigma$ for Gaussian noise; quality-loss units for JPEG (100 $-$ quality); radius for blur; px for letter spacing; class-distance units for font dimensions.}
\label{tab:phase1-thresholds}
\resizebox{\textwidth}{!}{%
\begin{tabular}{l cc cc cc}
\toprule
 & \multicolumn{2}{c}{Qwen2.5-VL} & \multicolumn{2}{c}{Gemini 3.1 Pro} & \multicolumn{2}{c}{GPT-4o} \\
\cmidrule(lr){2-3} \cmidrule(lr){4-5} \cmidrule(lr){6-7}
Dimension & 50\% tier & 50\% val. & 50\% tier & 50\% val. & 50\% tier & 50\% val. \\
\midrule
\texttt{blur}              & medium  & 3.22 rad          & \textbf{small} & \textbf{1.34 rad}   & \textbf{small} & \textbf{1.34 rad} \\
\texttt{gaussian\_noise}   & medium  & $\sigma$ 16.8     & medium         & $\sigma$ 16.8       & \textbf{small} & \textbf{$\sigma$ 4.7} \\
\texttt{color\_offset}     & xlarge  & $\Delta E$ 43.4   & medium         & $\Delta E$ 10.6     & strong         & $\Delta E$ 15.4 \\
\texttt{position\_offset}  & NC      & ---               & large          & 100.8 px            & NC             & --- \\
\texttt{rotation}          & xlarge  & 23.3$^\circ$      & large          & 18.1$^\circ$        & large          & 18.1$^\circ$ \\
\texttt{scale\_error}      & xxlarge & 77.4\%            & large          & 25.0\%              & xxlarge        & 77.4\% \\
\texttt{letter\_spacing}   & NC      & ---               & medium         & 3.5 px              & large          & 11.4 px \\
\texttt{opacity}           & large   & 0.694             & medium         & 0.302               & large          & 0.694 \\
\texttt{font\_family}      & NC      & ---               & large          & 1.0                 & large          & 1.0 \\
\texttt{font\_style}       & NC      & ---               & medium         & 0.5                 & NC             & --- \\
\texttt{font\_weight\_heavy} & NC    & ---               & small          & 100 CSS             & NC             & --- \\
\texttt{font\_weight\_light} & NC    & ---               & small          & 100 CSS             & NC             & --- \\
\texttt{jpeg\_compression} & NC      & ---               & NC             & ---                 & medium         & quality $-49.7$ \\
\bottomrule
\end{tabular}%
}
\end{table}

Key threshold findings: for image-level degradations, GPT-4o and Gemini reach 50\% blur detection at the small tier (radius $\approx 1.34$), while Qwen requires medium blur (radius $\approx 3.22$). For spatial degradations, thresholds are substantially later than what human observers require: scale errors must reach 25--77\% deviation, and position offsets must exceed 100 pixels, before any model reliably detects them. Colour thresholds range from $\Delta E$ 10.6 (Gemini) to $\Delta E$ 43.4 (Qwen), with the Qwen threshold near the ceiling of the $\Delta E$ scale used in this study. For typographic dimensions, only Gemini has defined thresholds (font weight: small tier; font style: medium tier); GPT-4o and Qwen do not cross 50\% on these dimensions.

The practical implication is that for the majority of error types, current VLM judges operate at perceptual thresholds that are far above the error magnitudes likely encountered in real text-editing evaluation. This perceptual limitation is compounded in Experiment 2 as the quality-discrimination thresholds are generally at least as late as the perceptual thresholds, and often later.

\subsection{Correlation with Phase 2 Patterns} \label{sec:exp1-exp2-correlation}

The two-experiment design allows a direct comparison between Experiment 1 (perceptual similarity score) and Experiment 2 (edit judgment score) on matched stimuli. The cross-experiment Spearman correlation quantifies global coherence between the two assessment modes.

\begin{table}[htbp]
\centering
\caption{Global Spearman correlations between Experiment 1 similarity score and Experiment 2 quality metrics on matched stimuli. $\rho(\text{sim} \leftrightarrow \text{OQ})$ = Spearman between Exp1 similarity and Exp2 \texttt{overall\_quality}, within-model rank analysis. $\rho(\text{sim} \leftrightarrow \text{NM})$ = Spearman between Exp1 similarity and Exp2 \texttt{narrow\_mean}, cross-model comparable.}
\label{tab:exp1-exp2-corr}
\begin{tabular}{lccc}
\toprule
Model & Matched stimuli & $\rho(\text{sim} \leftrightarrow \text{OQ})$ & $\rho(\text{sim} \leftrightarrow \text{NM})$ \\
\midrule
Qwen2.5-VL              & 623 & \textbf{0.301} & \textbf{0.224} \\
Gemini 3.1 Pro Preview  & 568 & \textbf{0.738} & \textbf{0.755} \\
GPT-4o                  & 621 & \textbf{0.492} & \textbf{0.509} \\
\bottomrule
\end{tabular}
\end{table}

All correlations are positive and significant ($p < 0.001$). Gemini shows a strong global correlation ($\rho(\text{sim} \leftrightarrow \text{NM}) = 0.755$): its perceptual judgment and edit quality judgment are tightly coupled. GPT-4o shows a moderate correlation ($\rho = 0.509$). Qwen shows the weakest correlation ($\rho(\text{sim} \leftrightarrow \text{NM}) = 0.224$), with the \texttt{narrow\_mean} correlation weaker than the OQ correlation.

The two-experiment design predicts that Experiment 1 perceptual blindness should transfer to Experiment 2 judgment failure: a model that cannot see a degradation cannot penalise it.

For Qwen, the prediction holds on three of four blind cells (\texttt{font\_style}, \texttt{font\_weight\_heavy}, \texttt{position\_offset}), all of which sit at or above the 4.10 baseline. \texttt{font\_weight\_light} is the exception (NM 3.94, ``always below baseline''), but given Qwen's broader pattern of noisy and inconsistent judgments across the rubric, this single deviation is more plausibly attributed to judgment noise than to a meaningful compensatory mechanism.

The Gemini case is more interesting. \texttt{jpeg\_compression} is perceptually blind in Experiment 1 (8.7\%, n/s) but produces a significant magnitude correlation in Experiment 2 ($\rho = -0.490^{**}$) with a quality-discrimination threshold at the medium tier. One plausible explanation is training-data exposure: ``spot-the-difference'' or image-quality datasets might pair lossless originals with JPEG versions labelled as visually equivalent or near-equivalent. In other words, the model can see the JPEG artefacts but has learned to ignore them when asked about similarity. Experiment 2 asks a different question, so they surface.

GPT-4o has no Experiment 1 blind cells under the binomial criterion, so the prediction is not directly testable; its weakest perceptual dimensions are also its weakest in Experiment 2, consistent with the prediction in trend.

The supported claim is the asymmetric one: Experiment 1 perceptual blindness usually transfers to Experiment 2, and Experiment 2 rarely recovers signal that Experiment 1 misses. Experiment 1 is a useful but not infallible upper bound on Experiment 2 sensitivity.
